% ─────────────────────────────────────────────────────────────────────────────
% 3D Gaussian Splatting – Beamer prezentácia pre cvičenie PVSO
% Kompilovať: pdflatex prezentacia.tex  (alebo Overleaf)
% ─────────────────────────────────────────────────────────────────────────────
\documentclass[aspectratio=169]{beamer}

% ── Moderný čistý vzhľad ─────────────────────────────────────────────────────
\usetheme{default}
\setbeamertemplate{navigation symbols}{}
\usefonttheme[onlymath]{serif}
\usefonttheme{professionalfonts}

% ── Farby ────────────────────────────────────────────────────────────────────
\definecolor{pvsoBlue}{RGB}{0,70,140}
\definecolor{pvsoRed}{RGB}{180,0,0}
\definecolor{pvsoGray}{RGB}{90,90,90}
\definecolor{pvsoLight}{RGB}{232,240,252}
\definecolor{pvsoDark}{RGB}{15,35,65}

% ── Základné farby tém ────────────────────────────────────────────────────────
\setbeamercolor{background canvas}{bg=white}
\setbeamercolor{normal text}{fg=pvsoDark}
\setbeamercolor{alerted text}{fg=pvsoRed}
\setbeamercolor{structure}{fg=pvsoBlue}

% ── Odrážky ──────────────────────────────────────────────────────────────────
\setbeamertemplate{itemize item}{%
  \color{pvsoBlue}\raisebox{0.5pt}{\footnotesize$\blacktriangleright$}}
\setbeamertemplate{itemize subitem}{%
  \color{pvsoBlue!65}\raisebox{0.5pt}{\scriptsize$\triangleright$}}
\setbeamertemplate{enumerate item}{%
  \textcolor{pvsoBlue}{\bfseries\insertenumlabel.}}
\setbeamertemplate{enumerate subitem}{%
  \textcolor{pvsoBlue!65}{\bfseries\insertsubenumlabel.}}

% ── Bloky – moderný plochý dizajn ────────────────────────────────────────────
\setbeamercolor{block title}{bg=pvsoBlue,fg=white}
\setbeamercolor{block body}{bg=pvsoLight,fg=pvsoDark}
\setbeamercolor{block title alerted}{bg=pvsoRed,fg=white}
\setbeamercolor{block body alerted}{bg=pvsoRed!10,fg=pvsoDark}
\setbeamertemplate{blocks}[rounded][shadow=false]

% ── Nadpis snímky ────────────────────────────────────────────────────────────
\setbeamercolor{frametitle}{fg=pvsoBlue,bg=white}
\setbeamerfont{frametitle}{size=\large,series=\bfseries}
\setbeamertemplate{frametitle}{%
  \vspace*{1pt}%
  {\usebeamerfont{frametitle}\usebeamercolor[fg]{frametitle}\insertframetitle}%
  \vspace*{1pt}%
  {\color{pvsoBlue}\hrule height 1.2pt}%
  \vspace*{2pt}%
}

% ── Pätička s progresovou lištou ─────────────────────────────────────────────
\setbeamercolor{progress done}{bg=pvsoBlue}
\setbeamercolor{progress left}{bg=pvsoBlue!20}
\setbeamertemplate{footline}{%
  \hbox{%
    \begin{beamercolorbox}[wd=\dimexpr\paperwidth*\insertframenumber/%
        \inserttotalframenumber\relax, ht=2.5pt, dp=0pt]{progress done}%
    \end{beamercolorbox}%
    \begin{beamercolorbox}[wd=\dimexpr\paperwidth-%
        \paperwidth*\insertframenumber/\inserttotalframenumber\relax,%
        ht=2.5pt, dp=0pt]{progress left}%
    \end{beamercolorbox}%
  }%
  \nointerlineskip%
  \begin{beamercolorbox}[wd=\paperwidth, ht=2.4ex, dp=1.2ex]{footline}%
    {\hspace{5pt}\tiny\color{pvsoGray}%
      \insertshortauthor\enspace–\enspace\insertshorttitle}%
    \hfill%
    {\tiny\color{pvsoGray}%
      \insertframenumber\,/\,\inserttotalframenumber\hspace{5pt}}%
  \end{beamercolorbox}%
}

% ── Titulná strana ────────────────────────────────────────────────────────────
\setbeamertemplate{title page}{%
  \begin{tikzpicture}[remember picture, overlay]
    \fill[pvsoBlue]
      (current page.north west) rectangle
      ([yshift=-0.54\paperheight]current page.north east);
    \node[anchor=center, text=white, font=\Large\bfseries,
          text width=0.82\paperwidth, align=center]
      at ([yshift=0.115\paperheight]current page.center) {\inserttitle};
    \node[anchor=center, text=white!80!pvsoBlue, font=\normalsize,
          text width=0.82\paperwidth, align=center]
      at ([yshift=0.022\paperheight]current page.center) {\insertsubtitle};
    \draw[white!45, line width=0.5pt]
      ([xshift=-0.28\paperwidth, yshift=-0.028\paperheight]current page.center) --
      ([xshift= 0.28\paperwidth, yshift=-0.028\paperheight]current page.center);
    \node[anchor=center, text=pvsoGray, font=\normalsize]
      at ([yshift=-0.13\paperheight]current page.center) {\insertauthor};
    \node[anchor=center, text=pvsoGray, font=\small]
      at ([yshift=-0.19\paperheight]current page.center) {\insertinstitute};
    \node[anchor=center, text=pvsoGray, font=\small]
      at ([yshift=-0.25\paperheight]current page.center) {\insertdate};
  \end{tikzpicture}%
}

% ── Balíčky ──────────────────────────────────────────────────────────────────
\usepackage[utf8]{inputenc}
\usepackage[T1]{fontenc}
\usepackage[slovak]{babel}
\usepackage{graphicx}
\usepackage{tikz}
\usetikzlibrary{arrows.meta, shapes, positioning, fit, backgrounds, calc}
\usepackage{amsmath}
\usepackage{booktabs}
\usepackage{fontawesome5}

% ── Metadáta ─────────────────────────────────────────────────────────────────
\title[3D Gaussian Splatting]{3D Gaussian Splatting}
\subtitle{Od fotografií k interaktívnej 3D scéne}
\author{PVSO – Cvičenie}
\institute{Počítačové videnie a spracovanie obrazu}
\date{2026}

% ── Pomocné príkazy ───────────────────────────────────────────────────────────
\newcommand{\highlight}[1]{\textcolor{pvsoBlue}{\textbf{#1}}}
\newcommand{\red}[1]{\textcolor{pvsoRed}{\textbf{#1}}}

% =============================================================================
\begin{document}

% ── Titulný slajd ─────────────────────────────────────────────────────────────
\begin{frame}[plain]
\titlepage
\end{frame}

% =============================================================================
% SLAJD 1 – Pipeline prehľad
% =============================================================================
\begin{frame}{Čo dnes urobíme}

\begin{tikzpicture}[
    node distance=0.4cm and 0.3cm,
    box/.style={rectangle, rounded corners=5pt, minimum width=2.6cm,
                minimum height=1cm, align=center, font=\small\bfseries},
    arrow/.style={-{Stealth[length=6pt]}, thick, pvsoBlue}
]
\node[box, fill=pvsoBlue!12, draw=pvsoBlue] (cam)
    {\faCamera\\[2pt] Fotky\\scény};
\node[box, fill=pvsoBlue!12, draw=pvsoBlue, right=of cam] (feat)
    {\faSitemap\\[2pt] COLMAP\\príznaky};
\node[box, fill=pvsoBlue!12, draw=pvsoBlue, right=of feat] (rec)
    {\faCube\\[2pt] COLMAP\\rekonštrukcia};
\node[box, fill=pvsoRed!12, draw=pvsoRed, right=of rec] (train)
    {\faBrain\\[2pt] Gaussian\\trénovanie};
\node[box, fill=green!12, draw=green!55!black, right=of train] (view)
    {\faEye\\[2pt] SIBR\\Viewer};

\draw[arrow] (cam) -- (feat);
\draw[arrow] (feat) -- (rec);
\draw[arrow] (rec) -- (train);
\draw[arrow] (train) -- (view);

\node[font=\tiny, pvsoGray, below=0.15cm of feat]  {5–30 min};
\node[font=\tiny, pvsoGray, below=0.15cm of rec]   {5–30 min};
\node[font=\tiny, pvsoGray, below=0.15cm of train] {30–60 min};
\node[font=\tiny, pvsoGray, below=0.15cm of view]  {real-time};
\end{tikzpicture}

\vspace{0.45cm}
\begin{columns}[t]
\column{0.5\textwidth}
\textbf{Teória (táto prezentácia)}
\begin{itemize}\setlength{\itemsep}{2pt}
    \item Ako COLMAP „vidí" 3D zo 2D
    \item Čo sú 3D Gaussiány
    \item Ako prebieha trénovanie
\end{itemize}

\column{0.5\textwidth}
\textbf{Cvičenie}
\begin{itemize}\setlength{\itemsep}{2pt}
    \item Nasnímať scénu kamerou Ximea
    \item Spustiť COLMAP GUI v Dockeri
    \item Natrénovať vlastný model
\end{itemize}
\end{columns}
\end{frame}

% =============================================================================
% SLAJD 2 – Problém: 3D z 2D
% =============================================================================
\begin{frame}{Problém: ako dostať 3D z 2D fotografií?}

\begin{columns}[c]
\column{0.52\textwidth}

Fotografia je \red{strata informácie} – 3D scéna sa premietne do 2D:
\[
\mathbf{p}_{2D} = \mathbf{K} \cdot [\mathbf{R} \mid \mathbf{t}] \cdot \mathbf{P}_{3D}
\]

\vspace{0.08cm}
\textbf{Čo vieme z jednej fotky?}
\begin{itemize}\setlength{\itemsep}{2pt}
    \item farbu pixelov \checkmark
    \item smer lúča kamery \checkmark
    \item \red{hĺbku (vzdialenosť)?} $\boldsymbol{\times}$
\end{itemize}

\vspace{0.08cm}
\textbf{Riešenie:} viac fotografií z rôznych uhlov\\
$\Rightarrow$ \highlight{triangulácia}

\column{0.48\textwidth}
\begin{center}
\begin{tikzpicture}[scale=1.15]
\draw[pvsoBlue, thick, fill=pvsoBlue!15]
    (0,3.2) -- (0.72,2.35) -- (-0.72,2.35) -- cycle;
\draw[pvsoBlue, thick, fill=pvsoBlue!15]
    (3.5,3.2) -- (4.22,2.35) -- (2.78,2.35) -- cycle;
\node[font=\small, pvsoBlue] at (0,3.55) {\textbf{Kamera 1}};
\node[font=\small, pvsoBlue] at (3.5,3.55) {\textbf{Kamera 2}};
\draw[pvsoRed, dashed, thick] (0,2.95) -- (1.75,1.05);
\draw[pvsoRed, dashed, thick] (3.5,2.95) -- (1.75,1.05);
\node[circle, fill=pvsoRed, inner sep=4pt] at (1.75,1.05) {};
\node[font=\small, pvsoGray, below=0.1cm] at (1.75,1.05) {3D bod (P)};
\node[font=\footnotesize, pvsoGray] at (1.75,0.32) {triangulácia};
\end{tikzpicture}
\end{center}
\end{columns}

\vspace{0.15cm}
\begin{alertblock}{Väzba na predchádzajúce prednášky}
{\small Kamera model $\mathbf{K}$, $[\mathbf{R}|\mathbf{t}]$ – PVSO prednáška 1.\enspace
Príznaky, Harris detektor, RANSAC – prednáška 7.}
\end{alertblock}

\end{frame}

% =============================================================================
% SLAJD 3 – COLMAP: príznaky a matching
% =============================================================================
\begin{frame}{COLMAP – Krok 1: Príznaky a Matching}

% Vizualizácia – centrovane, väčšia
\begin{center}
\begin{tikzpicture}[scale=1.0]
% Fotka 1
\draw[pvsoBlue, thick, fill=pvsoBlue!8] (0,0) rectangle (2.4,1.65);
\node[font=\small, pvsoBlue] at (1.2,1.95) {\textbf{Fotka 1}};
\node[circle, fill=pvsoRed, inner sep=3pt] at (0.55,1.15) {};
\node[circle, fill=pvsoRed, inner sep=3pt] at (1.25,0.60) {};
\node[circle, fill=pvsoRed, inner sep=3pt] at (1.95,1.30) {};

% Fotka 2
\draw[pvsoBlue, thick, fill=pvsoBlue!8] (3.1,0) rectangle (5.5,1.65);
\node[font=\small, pvsoBlue] at (4.3,1.95) {\textbf{Fotka 2}};
\node[circle, fill=pvsoRed, inner sep=3pt] at (3.40,1.20) {};
\node[circle, fill=pvsoRed, inner sep=3pt] at (4.25,0.55) {};
\node[circle, fill=pvsoRed, inner sep=3pt] at (5.10,1.35) {};

% Matching línie
\draw[green!60!black, thick] (0.55,1.15) -- (3.40,1.20);
\draw[green!60!black, thick] (1.25,0.60) -- (4.25,0.55);
\draw[green!60!black, thick] (1.95,1.30) -- (5.10,1.35);

% Popis
\node[font=\small, pvsoGray] at (2.75,-0.28) {korešpondencie (SIFT matching)};

% Šípka a výsledok
\node[font=\small, pvsoGray] at (2.75,-0.70) {$\Downarrow$\enspace triangulácia\enspace$\Rightarrow$\enspace sparse 3D body};

% Sparse cloud – kompaktne
\foreach \x/\y in {1.35/-1.35, 1.85/-1.10, 2.35/-1.40,
                   2.80/-1.13, 3.25/-1.33, 2.05/-1.55, 2.90/-1.53, 2.55/-1.23} {
    \node[circle, fill=pvsoBlue!75, inner sep=2.5pt] at (\x,\y) {};
}
\end{tikzpicture}
\end{center}

\vspace{0.04cm}
\begin{columns}[t]
\column{0.5\textwidth}
\begin{block}{Feature Extraction (SIFT)}
{\small Nájdenie \highlight{kľúčových bodov} – podobné Harrisovmu detektoru rohov (prednáška 7).
Pre každý bod: \textit{deskriptor} (128-dim vektor).}
\end{block}

\column{0.5\textwidth}
\begin{block}{Feature Matching}
{\small Nájdenie \highlight{korešpondenčných párov} medzi fotkami – porovnanie deskriptorov,
filtrácia pomocou \textbf{RANSAC} (prednáška 7).}
\end{block}
\end{columns}

\end{frame}

% =============================================================================
% SLAJD 4 – COLMAP: Rekonštrukcia
% =============================================================================
\begin{frame}{COLMAP – Krok 2: Sparse Rekonštrukcia a Undistortion}

\begin{columns}[t]
\column{0.50\textwidth}

\begin{block}{Sparse Reconstruction (SfM)}
Z korešpondencií COLMAP vypočíta:
\begin{itemize}\setlength{\itemsep}{2pt}
    \item \textbf{pozície kamier} $[\mathbf{R}|\mathbf{t}]$ v 3D priestore
    \item \textbf{sparse mrak bodov} – tisíce 3D bodov
\end{itemize}
{\footnotesize Výsledok: \texttt{cameras.bin}, \texttt{images.bin}, \texttt{points3D.bin}}
\end{block}

\vspace{0.08cm}
\begin{block}{Undistortion (\texttt{convert.py})}
{\small Reálne objektívy \textbf{skresľujú} obraz (radiálna distorzia). 3DGS potrebuje \textbf{ideálny pinhole model}:
\begin{itemize}\setlength{\itemsep}{0pt}
    \item COLMAP vypočíta parametre distorzie
    \item \texttt{convert.py} vytvorí undistortované fotky
\end{itemize}}
\end{block}

\column{0.50\textwidth}
\begin{center}
\begin{tikzpicture}[scale=1.1]
\fill[pvsoGray!15] (-2,-0.3) rectangle (2,0);
\foreach \x/\y in {-1.2/0.5, -0.5/1.0, 0/0.8, 0.5/1.3, 1.0/0.6,
                    1.4/0.9, -0.8/1.5, 0.3/0.4, -0.2/1.8, 0.8/0.2} {
    \node[circle, fill=pvsoBlue!75, inner sep=2pt] at (\x,\y) {};
}
\foreach \x in {-1.7, -0.85, 0, 0.85, 1.7} {
    \draw[pvsoRed, thick, fill=pvsoRed!20]
        (\x,2.55) -- (\x-0.27,2.12) -- (\x+0.27,2.12) -- cycle;
    \draw[pvsoRed, very thin] (\x,2.25) -- (0,0.9);
}
\node[font=\small, pvsoBlue] at (0,-0.62) {sparse point cloud};
\node[font=\small, pvsoRed]  at (0,2.92) {pozície kamier};
\end{tikzpicture}
\end{center}
\end{columns}

\end{frame}

% =============================================================================
% SLAJD 5 – Čo je 3D Gaussian Splatting?
% =============================================================================
\begin{frame}{3D Gaussian Splatting – základná myšlienka}

\begin{columns}[c]
\column{0.56\textwidth}

\textbf{Kerbl et al., SIGGRAPH 2023} –
scéna ako \highlight{milióny 3D Gaussiánov},
každý popisuje jeden „kúsok" scény.

\vspace{0.06cm}
\begin{block}{Prečo Gaussiány?}
\begin{itemize}\setlength{\itemsep}{1pt}
    \item \textbf{explicitná} reprezentácia (na rozdiel od NeRF)
    \item rýchle renderovanie – \textbf{real-time}
    \item ľahko diferencovateľné $\Rightarrow$ optimalizácia
    \item plynulé prechody a transparentnosť
\end{itemize}
\end{block}


\column{0.44\textwidth}
\begin{center}
\begin{tikzpicture}[scale=1.3]
\draw[pvsoBlue!85, fill=pvsoBlue!25, thick, rotate around={30:(1.1,2.7)}]
    (1.1,2.7) ellipse (0.88 and 0.40);
\draw[pvsoRed!85, fill=pvsoRed!25, thick, rotate around={-20:(2.85,2.85)}]
    (2.85,2.85) ellipse (0.66 and 0.30);
\draw[green!60!black, fill=green!22, thick, rotate around={65:(2.05,1.80)}]
    (2.05,1.80) ellipse (0.76 and 0.28);
\draw[orange!85!black, fill=orange!22, thick, rotate around={10:(0.90,1.20)}]
    (0.90,1.20) ellipse (0.58 and 0.37);
\draw[purple!80, fill=purple!22, thick, rotate around={-40:(3.15,1.85)}]
    (3.15,1.85) ellipse (0.70 and 0.28);
\draw[cyan!70!black, fill=cyan!18, thick, rotate around={25:(1.90,0.65)}]
    (1.90,0.65) ellipse (0.60 and 0.25);
\node[font=\small, pvsoGray] at (2.05,0.08) {milióny Gaussiánov = scéna};
\end{tikzpicture}
\end{center}
\end{columns}

\end{frame}

% =============================================================================
% SLAJD 6 – Parametre 3D Gaussiánu
% =============================================================================
\begin{frame}{3D Gaussián – čo ho popisuje}

Každý Gaussián má \highlight{tieto parametre} (optimalizujú sa pri trénovaní):

\vspace{0.05cm}
{\setlength\abovedisplayskip{2pt}\setlength\belowdisplayskip{2pt}%
\begin{columns}[t]
\column{0.25\textwidth}
\begin{block}{\faMapMarker\ Poloha}
\[\boldsymbol{\mu} \in \mathbb{R}^3\]
stred Gaussiánu v 3D priestore
\end{block}

\column{0.25\textwidth}
\begin{block}{\faCube\ Tvar}
\[\boldsymbol{\Sigma} \in \mathbb{R}^{3\times3}\]
kovariančná matica\\(veľkosť, orientácia)
\end{block}

\column{0.25\textwidth}
\begin{block}{\faEyeSlash\ Priehľadnosť}
\[\alpha \in [0,1]\]
opacity – ako veľmi je Gaussián „viditeľný"
\end{block}

\column{0.25\textwidth}
\begin{block}{\faPalette\ Farba}
Spherical Harmonics (SH)\\[2pt]
farba závisí od\\smeru pohľadu
\end{block}
\end{columns}}

\vspace{0.05cm}
\begin{center}
\begin{tikzpicture}[scale=0.95]
% ── Gaussián – elipsoid (reprezentuje Σ) ────────────────────────────────
\draw[pvsoBlue, thick, fill=pvsoBlue!18, fill opacity=0.65,
      rotate around={35:(0,0)}]
    (0,0) ellipse (1.85 and 0.72);

% Poloosi – naznačujú kovariančnú maticu
\draw[pvsoGray!65, dashed, rotate around={35:(0,0)}] (-1.85,0) -- (1.85,0);
\draw[pvsoGray!65, dashed, rotate around={35:(0,0)}] (0,-0.72) -- (0,0.72);

% ── Stred – μ ────────────────────────────────────────────────────────────
\node[circle, fill=pvsoRed, inner sep=4pt] (C) at (0,0) {};
\draw[-{Stealth[length=5pt]}, pvsoRed, thick] (0,0) -- (-0.5,1.0);
\node[pvsoRed, font=\small\bfseries, anchor=south] at (-0.5,1.02)
    {$\boldsymbol{\mu}$ – poloha stredu};

% ── Σ – tvar a orientácia (callout na obvod elipsy) ─────────────────────
\draw[pvsoBlue!80!black, thin, -{Stealth[length=4pt]}]
    (-3.1,0.1) -- (-1.62,0.35);
\node[pvsoBlue!80!black, font=\small\bfseries, anchor=east] at (-3.1,0.1)
    {$\boldsymbol{\Sigma}$ – tvar a orientácia};

% ── α – priehľadnosť (callout na výplň) ─────────────────────────────────
\draw[pvsoGray, thin, -{Stealth[length=4pt]}]
    (2.85,0.2) -- (1.1,0.1);
\node[pvsoGray, font=\small, anchor=west] at (2.85,0.2)
    {$\alpha$ – priehľadnosť};

% ── Farba – view-dependent: dve šípky pohľadov ──────────────────────────
\draw[-{Stealth[length=6pt]}, orange!80!red, thick]
    (-2.5,-1.0) -- (-0.55,-0.2);
\draw[-{Stealth[length=6pt]}, cyan!65!blue, thick]
    ( 2.5,-1.0) -- ( 0.6,-0.18);
\node[orange!80!red, font=\scriptsize, anchor=north] at (-2.3,-1.0) {pohľad 1};
\node[cyan!65!blue,  font=\scriptsize, anchor=north] at ( 2.3,-1.0) {pohľad 2};
\node[pvsoGray, font=\footnotesize] at (0,-1.35)
    {farba $= f(\text{smer pohľadu})$ – Spherical Harmonics};
\end{tikzpicture}
\end{center}

\end{frame}

% =============================================================================
% SLAJD 7 – Splatting: renderovanie
% =============================================================================
\begin{frame}{Splatting – ako sa 3D Gaussiány premietajú na 2D}

\begin{columns}[c]
\column{0.54\textwidth}

\textbf{Renderovanie} = premietnutie 3D Gaussiánov do 2D obrazu

\vspace{0.2cm}
\begin{enumerate}\setlength{\itemsep}{5pt}
    \item \highlight{Projekcia}: každý 3D Gaussián $\to$ 2D elipsa\\
    {\small(kamerová projekcia – prednáška 1)}

    \item \highlight{Zoradenie}: podľa hĺbky od najbližšieho

    \item \highlight{Alpha blending}: farby cez seba ($\alpha$):
    \[
    C = \sum_{i} c_i \cdot \alpha_i \prod_{j<i}(1-\alpha_j)
    \]
\end{enumerate}

\vspace{0.1cm}
Celý proces je \highlight{diferencovateľný}
$\Rightarrow$ môžeme optimalizovať parametre!

\column{0.46\textwidth}
\begin{center}
\begin{tikzpicture}[scale=1.12]
\node[font=\small, pvsoBlue] at (1.5,4.2) {\textbf{3D priestor}};
% 3D Gaussiány
\draw[pvsoBlue, fill=pvsoBlue!25, rotate around={30:(0.8,3.1)}]
    (0.8,3.1) ellipse (0.65 and 0.27);
\draw[pvsoRed, fill=pvsoRed!25, rotate around={-15:(2.3,3.3)}]
    (2.3,3.3) ellipse (0.55 and 0.23);
\draw[green!60!black, fill=green!25, rotate around={50:(1.6,2.55)}]
    (1.6,2.55) ellipse (0.60 and 0.23);
% Kamera
\draw[pvsoGray, thick, fill=pvsoGray!25]
    (1.5,1.95) -- (1.2,1.58) -- (1.8,1.58) -- cycle;
\node[font=\small, pvsoGray] at (1.5,1.33) {kamera};
% Premietacie lúče
\draw[pvsoGray, dashed] (0.8,2.92) -- (0.72,0.62);
\draw[pvsoGray, dashed] (2.3,3.10) -- (2.18,0.62);
% 2D obraz
\draw[thick, pvsoBlue] (-0.1,0.05) rectangle (3.1,0.68);
\node[font=\small, pvsoBlue] at (1.5,-0.28) {\textbf{2D obraz (pixel buffer)}};
% 2D elipsy
\draw[pvsoBlue, fill=pvsoBlue!35, rotate around={10:(0.72,0.36)}]
    (0.72,0.36) ellipse (0.50 and 0.17);
\draw[pvsoRed, fill=pvsoRed!35, rotate around={-5:(2.20,0.36)}]
    (2.20,0.36) ellipse (0.45 and 0.17);
\draw[green!60!black, fill=green!35, rotate around={15:(1.55,0.36)}]
    (1.55,0.36) ellipse (0.55 and 0.17);
\end{tikzpicture}
\end{center}
\end{columns}

\end{frame}

% =============================================================================
% SLAJD 8 – Trénovanie
% =============================================================================
\begin{frame}{Trénovanie – ako sa Gaussiány učia}

\begin{center}
\begin{tikzpicture}[
    node distance=0.5cm and 0.55cm,
    box/.style={rectangle, rounded corners=4pt, minimum width=2.65cm,
                minimum height=0.88cm, align=center, font=\small},
    arrow/.style={-{Stealth[length=5pt]}, thick}
]
\node[box, fill=pvsoBlue!12, draw=pvsoBlue] (init)
    {Inicializácia\\{\footnotesize sparse cloud z COLMAP}};
\node[box, fill=pvsoBlue!12, draw=pvsoBlue, right=1.1cm of init] (render)
    {Renderovanie\\{\footnotesize splatting $\to$ 2D obraz}};
\node[box, fill=pvsoRed!12, draw=pvsoRed, right=1.1cm of render] (loss)
    {Strata (Loss)\\{\footnotesize render vs. fotka}};
\node[box, fill=green!12, draw=green!55!black, below=0.9cm of loss] (optim)
    {Optimalizácia\\{\footnotesize $\mu, \Sigma, \alpha$, farba}};
\node[box, fill=orange!12, draw=orange!70!black, below=0.9cm of init] (dens)
    {Adaptívna\\densifikácia};

\draw[arrow, pvsoBlue] (init)   -- (render);
\draw[arrow, pvsoBlue] (render) -- (loss);
\draw[arrow, pvsoRed]  (loss)   -- (optim)
    node[midway, right, font=\footnotesize]{backprop};
\draw[arrow, green!55!black]    (optim) -- (dens);
\draw[arrow, orange!70!black]   (dens)  -- (init);

\node[font=\footnotesize, pvsoGray, below=0.12cm] at ($(dens)!0.5!(optim)$)
    {pridaj Gaussiány tam kde chýbajú};
\node[font=\footnotesize, pvsoGray, below=0.12cm of optim]
    {Adam optimizer};
\end{tikzpicture}
\end{center}

\vspace{0.2cm}
\begin{columns}[t]
\column{0.5\textwidth}
\textbf{Loss funkcia:}
\[
\mathcal{L} = (1-\lambda)\mathcal{L}_1 + \lambda \mathcal{L}_{D\text{-}SSIM}
\]
{\small porovnáva renderovaný obraz s originálom}

\column{0.5\textwidth}
\textbf{Priebeh trénovania:}
\begin{itemize}\setlength{\itemsep}{2pt}
    \item 30\,000 iterácií (default)
    \item checkpoint po 7\,000 a 30\,000
    \item výsledok: \texttt{point\_cloud.ply}
\end{itemize}
\end{columns}

\end{frame}

% =============================================================================
% SLAJD 9 – Výsledky
% =============================================================================
\begin{frame}{Výsledky a možnosti zobrazenia}

\begin{columns}[t]
\column{0.50\textwidth}

\begin{block}{Čo dostaneme?}
\begin{itemize}\setlength{\itemsep}{2pt}
    \item súbor \texttt{point\_cloud.ply} – milióny 3D Gaussiánov
    \item renderovanie z \textbf{ľubovoľného uhlu}
    \item \textbf{real-time} interaktívny viewer
\end{itemize}
\end{block}

\vspace{0.1cm}
\begin{block}{SIBR Gaussian Viewer}
\begin{itemize}\setlength{\itemsep}{2pt}
    \item interaktívna 3D navigácia
    \item ovládanie: WASD + myš
    \item pohľad z pozícií trénovacích kamier
\end{itemize}
\end{block}

\vspace{0.1cm}
\textbf{Metrika kvality:}\enspace
{\small PSNR $\approx$ 27–30 dB = dobrý výsledok\enspace
{\footnotesize(30+\,dB = výborný)}}

\column{0.50\textwidth}

\begin{block}{Typický výsledok – truck dataset}
\begin{tabular}{ll}
\toprule
Iterácie & PSNR \\
\midrule
7\,000   & $\approx$ 24.7 dB \\
30\,000  & $\approx$ 27.5 dB \\
\bottomrule
\end{tabular}
\end{block}

\vspace{0.1cm}
\begin{block}{Kde sa používa?}
\begin{itemize}\setlength{\itemsep}{2pt}
    \item VR/AR – interaktívne scény
    \item robotika – mapovanie prostredia
    \item kultúrne dedičstvo – digitalizácia
    \item filmová produkcia
    \item \textbf{priemyselná inšpekcia}
\end{itemize}
\end{block}

\end{columns}

\end{frame}

% =============================================================================
% SLAJD 10 – Cvičenie
% =============================================================================
\begin{frame}{Cvičenie – čo budete robiť}

\begin{columns}[t]
\column{0.52\textwidth}

\begin{block}{\faDocker\ Prostredie – Docker}
{\small Všetko je v Docker image:
\begin{itemize}\setlength{\itemsep}{1pt}
    \item COLMAP 3.9.1 (GUI + CUDA)
    \item 3D Gaussian Splatting
    \item SIBR Viewer
\end{itemize}
Spustenie: \texttt{./run.sh}}
\end{block}

\vspace{0.1cm}
\begin{block}{\faCamera\ Kamera Ximea}
{\small
\begin{itemize}\setlength{\itemsep}{1pt}
    \item priemyselná kamera s vysokým rozlíšením
    \item nasnímajte scénu z \textbf{min. 50 uhlov}
    \item každý bod viditeľný \textbf{aspoň 3$\times$}
    \item uložte do \texttt{data/moja\_scena/input/}
\end{itemize}}
\end{block}

\column{0.48\textwidth}

\begin{block}{\faListOl\ Postup}
{\footnotesize
\begin{enumerate}\setlength{\itemsep}{0pt}
    \item \texttt{./run.sh} – spusti Docker
    \item \texttt{colmap gui} – rekonštrukcia
    \begin{itemize}\setlength{\itemsep}{0pt}
        \item Feature extraction (SIMPLE\_PINHOLE)
        \item Feature matching (Exhaustive)
        \item Reconstruction
        \item Export model $\to$ \texttt{distorted/sparse/0/}
    \end{itemize}
    \item \texttt{python3 convert.py -s data/scena}
    \item \texttt{python3 train.py -s data/scena}
    \item \texttt{SIBR\_gaussianViewer\_app -m output/<uuid>}
\end{enumerate}}
\end{block}

\begin{alertblock}{}
Podrobný postup: \textbf{NAVOD.md}
\end{alertblock}

\end{columns}

\end{frame}

% =============================================================================
% SLAJD 11 – Tipy pre snímanie
% =============================================================================
\begin{frame}{Tipy pre kvalitné snímanie scény}

\begin{columns}[t]
\column{0.50\textwidth}
\begin{block}{\faCheck\ Správne}
\begin{itemize}\setlength{\itemsep}{2pt}
    \item \textbf{50–200 fotografií} scény
    \item veľké \textbf{prekrytie} – každý bod z $\geq$3 uhlov
    \item pomalý pohyb okolo objektu
    \item rôzne výšky kamery
    \item \textbf{statická scéna} – nič sa nesmie hýbať
    \item dobré rovnomerné osvetlenie
\end{itemize}
\end{block}

\column{0.50\textwidth}
\begin{alertblock}{\faTimes\ Problémy}
\begin{itemize}\setlength{\itemsep}{2pt}
    \item \red{lesklé / zrkadlové povrchy}
    {\footnotesize(odlesk mení farbu podľa uhlu)}
    \item \red{príliš málo fotografií}
    {\footnotesize(COLMAP nenájde dostatok zhôd)}
    \item \red{pohyb v scéne}
    {\footnotesize(auto, ľudia, vietor)}
    \item \red{homogénne plochy}
    {\footnotesize(biela stena – žiadne príznaky)}
\end{itemize}
\end{alertblock}
\end{columns}

\vspace{0.2cm}
\begin{center}
\begin{tikzpicture}[scale=0.92]
\draw[pvsoGray, dashed, thick] (0,0) ellipse (3.2 and 1.0);
\node[circle, fill=pvsoBlue, inner sep=4pt] at (0,0) {};
\node[font=\small, pvsoBlue] at (0,0.38) {\textbf{objekt}};
\foreach \a in {0,45,...,315} {
    \draw[pvsoRed, -{Stealth[length=8pt]}, thick]
        ({3.2*cos(\a)},{1.0*sin(\a)}) -- ({2.45*cos(\a)},{0.77*sin(\a)});
    \draw[pvsoRed, fill=pvsoRed!25, thick]
        ({3.2*cos(\a)},{1.0*sin(\a)})
        -- ({3.2*cos(\a)+0.30*cos(\a-90)},{1.0*sin(\a)+0.30*sin(\a-90)})
        -- ({3.2*cos(\a)+0.30*cos(\a+90)},{1.0*sin(\a)+0.30*sin(\a+90)}) -- cycle;
}
\node[font=\small, pvsoGray, anchor=west] at (3.55, 0.28) {ideálna trajektória};
\node[font=\small, pvsoGray, anchor=west] at (3.55,-0.22) {= kruh okolo objektu};
\end{tikzpicture}
\end{center}

\end{frame}

% =============================================================================
% SLAJD 12 – Záver
% =============================================================================
\begin{frame}{Zhrnutie}

\begin{columns}[c]
\column{0.55\textwidth}
\begin{block}{Čo sme sa naučili}
\begin{itemize}\setlength{\itemsep}{3pt}
    \item \highlight{COLMAP} (SfM): z 2D fotografií vypočíta polohy kamier
    a sparse 3D mrak bodov\\
    {\footnotesize – využíva príznaky + RANSAC (prednáška 7)}
    \item \highlight{3D Gaussian Splatting}: scéna = milióny 3D Gaussiánov
    s parametrami $(\mu, \Sigma, \alpha, \text{farba})$
    \item \highlight{Splatting}: projekcia Gaussiánov do 2D + alpha blending
    \item \highlight{Trénovanie}: optimalizácia cez diferencovateľné
    renderovanie
\end{itemize}
\end{block}

\column{0.45\textwidth}
\begin{alertblock}{Na cvičení}
\begin{enumerate}\setlength{\itemsep}{3pt}
    \item Nasnímaj scénu kamerou Ximea
    \item COLMAP GUI (rekonštrukcia)
    \item \texttt{convert.py} (undistortion)
    \item \texttt{train.py} (30–60 min)
    \item SIBR Viewer (výsledok!)
\end{enumerate}
\end{alertblock}

\vspace{0.15cm}
\begin{block}{Zdroje}
\begin{itemize}\small\setlength{\itemsep}{2pt}
    \item Kerbl et al.\ SIGGRAPH 2023
    \item \url{https://learnopencv.com/3d-gaussian-splatting/}
    \item \texttt{NAVOD.md} v repozitári
\end{itemize}
\end{block}
\end{columns}

\end{frame}

% =============================================================================
\end{document}
